\BourbakiUnnumberedChapter{致读者}{致读者}

\begin{enumerate}
\def\labelenumi{\arabic{enumi}.}
\item
  《数学原理丛书》从数学的起点开始，并给出完整的证明。原则上，它不要求读者具备特定的数学知识，而只要求对数学推理有一定的熟悉程度以及一定的抽象思维能力。尽管如此，它尤其面向那些至少已掌握大学数学课程第一年或第二年内容的读者。
\item
  我们选择的阐述方法是公理化的，通常从一般到特殊。证明的要求对主题内容施加了严格固定的顺序。因此，某些考虑的用处不会立即为读者所明了，除非他已经具备相当广泛的数学知识。
\item
  本丛书分为若干卷，每卷又分为若干章。以下列出法文版中已全部或部分出版的各卷。当英文译本可用时，相应的英文书名在括号中给出。全书中，引用标注指出英文版（若可用），否则指出法文版。
\end{enumerate}

Théorie des Ensembles（集合论）记作 E（Set Theory）Algèbre（代数 \(^{(1)}\) ）--- A（Alg）Topologie Générale（一般拓扑学）--- TG（Gen.～Top.）Fonctions d'une Variable Réelle --- FVR Espaces Vectoriels Topologiques（拓扑向量空间）--- EVT（Top. Vect. Sp.）Intégration --- INT Algèbre Commutative（交换代数 \(^{(2)}\) ）--- AC（Comm. Alg.）Variétés Différentielles et Analytiques --- VAR Groupes et Algèbres de Lie（李群与李代数 \(^{(3)}\) ）--- LIE（Lie）Théories Spectrales --- TS

在前六卷（按上述顺序）中，正文中的每一陈述都假定只已知那些已在同一章或按如下顺序排列的前面各章中讨论过的结果：E；A，第I至III章；TG，第I至III章；A，自第IV章起；TG，自第IV章起；FVR；EVT；INT。

从第七卷起，读者通常会发现对其与其他各卷逻辑关系的精确说明（前六卷始终假定为已知）。

\begin{enumerate}
\def\labelenumi{\arabic{enumi}.}
\setcounter{enumi}{3}
\item
  然而，我们有时在正文中插入一些例子，这些例子涉及读者可能已经知道但尚未在本丛书中讨论过的事实。这类例子置于两个星号之间：*\ldots*。毫无疑问，大多数读者会发现这些例子有助于理解正文。在其他情况下，两个星号之间的段落 *\ldots* 涉及本丛书其他部分讨论的结果。我们希望读者能够验证不存在任何恶性循环。
\item
  每一章的逻辑框架由该章的定义、公理和定理组成。这些是后续使用中主要需要牢记的部分。较不重要的结果以及那些容易从定理推导出的结果被标注为“命题”、“引理”、“推论”、“注记”等。那些在初读时可以略过的内容以小号字体印刷。对特别重要定理的评注偶尔以“附释”之名出现。
\end{enumerate}

为避免冗长的重复，有时引入仅在某一章或某一章中某一节内有效的记号或缩写是方便的（例如，在仅涉及交换环的章中，“环”一词总是指“交换环”）。这类约定总是被明确提及，通常出现在它们所属章的开头。

\begin{enumerate}
\def\labelenumi{\arabic{enumi}.}
\setcounter{enumi}{5}
\tightlist
\item
  正文中的某些段落旨在预先警告读者避免严重错误。这些段落在页边以符号 \(\sum\) （“危险弯道”）标出。
\end{enumerate}

\pandocbounded{\includegraphics[keepaspectratio]{images/part001_8ff0da04.jpg}}

\begin{enumerate}
\def\labelenumi{\arabic{enumi}.}
\setcounter{enumi}{6}
\item
  习题的设计既使读者能够确信自己已消化正文，又使其注意到正文中没有位置但仍有意义的结果。最困难的习题带有符号Ⅱ。
\item
  总的来说，我们遵循了普遍接受的术语，除非在看来有充分理由偏离它们的情况下。
\item
  我们特别努力始终使用严格正确的语言，同时不牺牲简洁性。我们在正文中尽可能提请注意语言的滥用；没有这些滥用，任何数学文本都有沦为迂腐、甚至不可读的风险。
\item
  由于原则上本文本是对一个理论的教条式阐述，它一般不含文献引用。参考文献集中在历史注记中。每个历史注记后的参考书目一般只包含在所讨论理论的演变中具有最重要意义的书籍和原始论文。它绝不自诩完备。
\end{enumerate}

至于习题，我们一般认为没有必要指出它们的来源，因为它们取自许多不同的资料（原始论文、教科书、习题集）。

\begin{enumerate}
\def\labelenumi{\arabic{enumi}.}
\setcounter{enumi}{10}
\tightlist
\item
  在本书中，对定理、公理、定义的引用， \ldots{} 通过依次引用以下内容给出：
\end{enumerate}

—— 册（使用第3节所列的缩写）、章和页，在那里可以找到它们；

—— 当涉及本书时，仅引用章和页。

结果概要用字母R引用；因此《集合论，R》表示“集合论结果概要”。

